\documentclass[11pt]{article}
\usepackage{amsmath,amssymb,amsthm,hyperref}
\newtheorem{theorem}{Theorem}\newtheorem{proposition}[theorem]{Proposition}
\title{Guarded incremental arena feedback and its native cost-trace compiler}
\date{16 November 2026}
\begin{document}\maketitle
\paragraph{Audit and result sought.} Namespace-refined compiler with written proofs and independent mathematical audits. Global identifier inverses are material: per-configuration bijections alone do not give a global inverse. The actual loop-labelled deployment and invocation paths are specified below. Exact finite diagnostics supplement the written proofs; no novelty, compilation, or formal-verification claim is made. The input is an independently defined incremental arena, not an atomic replacement of its complete execution. Its early loop outputs start child invocations immediately, before the parent invocation completes. The output is an actual locally finite recurrent arena with full invocation grade $R_UC$. Independent terminal semantics is related by native isomorphism to invocation-level subprobability recursion and to a countable cost trace. Shared event controls remain available and strictly exceed the invocation-atomic conditional regime. Mid-invocation prefixes are operational configurations, not fictitious known-$C$ partial trees.

\section{Input axioms and separate coefficients}
Use the locally finite, possibly unbounded incremental arena interface of the retained arena source. Explicitly, $P:A+U\rightsquigarrow B+U$ has all configurations $S_x$ in each input fibre; designated initials $i(x)$; finite enabled event sets; finite normalized raw rows at every enabled event with actual augmented context coefficients in $\mathrm{Vect}_k/k$; and deterministic transitions. Events execute at most once per invocation. Simultaneously enabled events persist with literally unchanged native rows and satisfy the actual diamond/cube equations. Each configuration has finitely many stable ordered emitted ports, whose labels and relative order persist; births are unique, and the port inclusions commute in the diamonds. New ports may be inserted between previously emitted ones. Empty enabled set means terminal. There is no bound on the length of all raw runs.

Independently, $P$ has one full weighted polynomial collection $C$, with coefficient $D_c$ in a symmetric monoidal content category $\mathcal D$, and a terminal shape $c(s)\in C_{|O(s)|}$ at every terminal configuration. No $C$ shape is required at a nonterminal configuration. A native arena map preserves all configurations, input anchors and initials, has bijections on enabled events and emitted ports at every compared configuration, actual exact-pushforward row maps coherent with transition/persistence, and one full arity-preserving weighted collection map satisfying the terminal body square. Configurations outside designated initial reach and all zero outcomes remain part of these data. Changing-endpoint maps use the same map $u:U\to U'$ on the incoming and outgoing loop tag, and preserve the $A/U$ and $B/U$ tags.

The guard is
\[
 \text{no initial configuration }i(x)\text{ emits a }U\text{-labelled port}. \tag{G}
\]
Initial $B$ ports are allowed, including a zero-event terminal invocation. Consequently every $U$ port is born at a paid primitive event, and initializing a child cannot recursively initialize further children without another paid event. This is a structural sufficient guard, not a claim that all useful feedback requires it.

Context coefficients belong to events; content coefficients belong independently to complete invocation shapes. No arrow from a tensor of event context coefficients to a content coefficient is supplied or needed. All context products in this manuscript are tensor products, with their augmentation. Countable effect/row aggregation uses the reviewed Vect and extended-nonnegative ambient. Finite histories have finite coefficient tensors. All full unused $C$ and $R_UC$ shapes are retained.

\section{The actual feedback configuration compiler}
For each input $x\in A+U$, define $\operatorname{Rec}_P(x)$ to be all finite rooted invocation trees with the following data. Each node $v$ carries a configuration $s_v$ of $P$ in its appropriate input fibre. For every emitted $U$ port $p$ of $s_v$ there is exactly one child invocation, keyed by the stable pair $(p,u_p)$ of its port identifier and actual loop label, with input $u_p$. There are no other children. Every such finite compatible tree is included, not only those reachable from the designated initial states. Unexecuted external roots retain their input tags; no external $A$ value is recast as a $U$ state.

This data type has the exact root/children description
\[
 \operatorname{Rec}_P(x)\cong
 \coprod_{s\in S_x}\prod_{p\in O_U(s)}\operatorname{Rec}_P(u_p),       \tag{1}
\]
restricted to well-founded finite trees. It is the initial algebra in set-valued families for the displayed finitary configuration functor: fold a finite tree by first folding its children and then applying the supplied algebra at its root. Structural induction proves existence and uniqueness and the algebra equation. This is a configuration-data universal property only, not a universal feedback completion or a claim that a probabilistic least fixed point has already been selected. The designated initial feedback configuration is the single node $i(x)$, which exists without silent descendants by (G).

The enabled events are all pairs $(v,e)$ with $e$ enabled in $s_v$. Its row is the original row at $(s_v,e)$. A transition updates that one configuration, keeps every existing child invocation, and initializes one child at each newly emitted $U$ port. No further child initialization is needed at this transition, by (G). Invocation identifiers are finite words of pairs $(p,u)$, retaining the actual loop label on every edge, not currently assigned integer output ranks. An event identifier is the typed pair of that word and its original local event identifier; an output-port identifier is the typed pair of the word and its original local port identifier. Root words are empty and are not tagged by the external input $a$. A port from an early composite branch can precede or follow later ports only once the final full order is known; its stable identifier and loop label are retained throughout. The same original port identifier can carry different labels in different configurations, so the label is not redundant global namespace data.

Feedback output ports are the $B$ ports found by recursively substituting each child invocation's emitted ports into its parent's corresponding $U$ port. Their identifiers are the stable invocation/port paths; their order is the induced recursive planar order. A transition can add new output ports or insert them between existing ones, but it cannot reverse the relative order of old ports. The requested feedback arena has external roots $a\in A$; the $x\in U$ configurations in (1) supply the recursively used continuation fibres.

\begin{theorem}[Native guarded feedback arena]
The compiler defines a locally finite recurrent incremental arena $\operatorname{Fb}_U(P):A\rightsquigarrow B$. Its full terminal collection is $R_UC$. It is functorial on all guarded native arena maps, retaining actual enabled-row arrows, unused configurations, and the entire induced content map $R_u\gamma$.
\end{theorem}
\begin{proof}
Each configuration has finitely many invocation nodes, each with finite enabled and emitted sets. A transition adds only finitely many children, and each new initial child has no $U$ port. Hence every successor is again a finite compatible configuration; enabled sets and emitted outputs remain finite. The event row is normalized because it is one original row. No event identifier can repeat: it is an original one-time event within its stable invocation path, and a child has a unique birth port.

For two enabled events in the same invocation, use the original diamond. Their distinct fresh port births initialize exactly the same children with the same input labels in either order; existing children are untouched. For events in different invocations, their local updates are disjoint even when one invocation is an ancestor of the other: updating the ancestor never changes an already emitted port's label or the existing descendant configuration. Fresh children cannot coincide with previously existing invocations. Thus the two updates commute. All surviving rows are literally unchanged and the port inclusions agree. The same componentwise reasoning gives every finite cube. The recursive output ordering respects port persistence at each node, so all old relative orders and labels remain unchanged.

Terminality means every invocation configuration is terminal. Each then supplies its full local shape $c(s_v)$ with its ordered ports. An emitted $U$ port corresponds to its child invocation, while a $B$ port is a settled leaf. Thus these local shapes render a finite $R_UC$ tree and exactly its ordered $B$ leaves. The coefficient at that full shape is $\bigotimes_vD_{c(s_v)}$. It is not inferred from event context coefficients. The global grade is declared to be all $R_UC$, including every tree with no realizing configuration.

If its original port function is $\omega$ and loop-label map is $u_*$, a native map sends every invocation edge $(p,u)$ to $(\omega(p),u_*(u))$, and applies the original global event/port functions to the final local identifier. This is a single global function on each occurrence union, even when $u_*$ is noninjective. Distinct simultaneously present ports remain distinct by the per-configuration port bijection. Pure external-input reindexing changes no identifier by an extra input tag, and pure $B$-output transport changes no invocation label. A native map acts on every invocation configuration and on all port/event identifiers. Port bijections give a bijection of the finite invocation nodes and preserve the appropriate incoming loop labels. Initial configurations and the new-child initialization rule commute. Transition squares and literal row persistence lift componentwise, giving the actual row maps and a bijection of enabled events. At terminal states the local full body squares assemble into the $R_u\gamma$ body square, using its actual content maps on every full tree. Identities and composites agree on each node, event and port and on all full grade factors. This proves functoriality without restricting to reachable configurations.
\end{proof}

This unfolding is a finite current configuration with potentially unbounded future behavior. It is not an invocation barrier: a child can execute while its parent has further enabled events. Neither the old finite-event-bound arena hypothesis nor an arbitrary trace axiom is being silently imposed.

\section{A native initial-algebra characterization of the loop fibre}
The pinned arena source uses the following explicit identifier convention: the identifier sets of an arena consist of identifiers occurring in enabled or emitted sets of some configuration, including unreachable configurations. A native identifier map has all those coordinates and no independent values on extraneous ambient namespace names. If additional phantom-name maps are included as independent morphism data, the uniqueness claim below need not determine them; that enlarged interface is not being quotiented implicitly.

Work in the intrinsic native hom category $\mathsf{Arena}(U,B)$. No strategy is selected. First define the loop-labelled deployment/case arena $\operatorname{Case}^{\rm lab}_U(X):B+U\rightsquigarrow B$ independently. At input $b\in B$ it is the event-free identity with one fixed branch-tagged singleton port identifier, independent of $b$. At input $u\in U$ it has exactly the $u$-fibre configurations of $X$, with event identifiers $(\mathrm{loop},u,e)$ and port identifiers $(\mathrm{loop},u,q)$. Its enabled rows and transitions are those of $X$. On a fixed-$U$ native map, its global identifier functions send $(u,e)$ to $(u,\eta(e))$ and $(u,q)$ to $(u,\omega(q))$; changing the loop interface sends the displayed $u$ to its supplied image as well. This is a tagged disjoint input-fibre construction, hence an actual arena and functor on $X$. It is not ordinary pure input reindexing, and no input-dependent tag is added to the structural identity branch.

Define $T_P(X)=P|_U;\operatorname{Case}^{\rm lab}_U(X)$, with the event-free identity copies removed by their structural configuration identification. At input $u$, a configuration of $T_P(X)$ is a root configuration $s\in S_u$ of $P$, together with one $X$ configuration at input $u_p$ for each emitted $U$ port $p$ of $s$. Continuation event and port identifiers are the typed triples $(p,u_p,e_X)$ and $(p,u_p,q_X)$; root-event and root-$B$-port identifiers use separate typed constructors. Its events are the root's enabled events and the enabled events in these copies. New $U$ ports initialize new $X$ copies immediately. Root $B$ ports are structural identity outputs, not sampled events. The output ports are root $B$ ports and the outputs of the $X$ copies in the root's ordered port slots. Forgetting the operational loop-label tags gives a directed native map to the unlabelled-deployment case with the same configuration data and corresponding full grade. It need not have a global identifier inverse if one $X$ identifier occurs over different loop inputs. Therefore this theorem concerns the explicitly labelled endofunctor, not an asserted operational isomorphism to an unmodified unlabelled case.

Let $G_X$ be the entire content collection of $X$. The full case grade is
\[
 I+U\cdot G_X,
\]
where $I$ is the unary unit collection and $U\cdot G_X$ has all pairs $(u,g)$ with coefficient that of $g$, even if that tag/shape never appears. A root $B$ port uses the $I$ shape; a root $U$ port labelled $u$ uses $(u,g)$ from its terminal copy. The full grade of $T_P(X)$ is
\[
 T_P^{\rm gr}(G_X)=C\circ(I+U\cdot G_X).                  \tag{4}
\]
The $U$ tags cannot be omitted: they are the internal loop labels of the eventual free tree grade. On native maps, $T_P$ is identity on the root $P$ data and uses the given map in every $X$ copy, with the full grade map induced by (4). The arena open-composition equations show this is a functor. In particular it acts on all configurations, not only on selected executions.

Write $F=\operatorname{Rec}_P|_U$ for the compiled feedback arena with roots in $U$. Root/children grafting gives an operational isomorphism $a:T_P(F)\to F$. Its global event map sends a continuation identifier $(p,u,(v,e))$ to $((p,u)\cdot v,e)$, where $v$ is the child invocation word; its inverse removes the first labelled pair. The same rule applies to output identifiers, with empty-word root cases handled by their separate constructors. These are genuine global inverse functions, not functions selected afresh at each configuration. Its full grade map
\[
 C\circ(I+U\cdot R_UC)\longrightarrow R_UC
\]
is the root constructor for nonempty finite trees, an isomorphism on all shapes with its actual tensor unitors and associators. A native $T_P$-algebra is an arena $X$ with a native arrow $\alpha:T_P(X)\to X$. This is an independently specified class: enabled/port bijections, exact row maps and all full coefficient/grade arrows are the ordinary arena map requirements, not assumed feedback-preservation equations. Algebra morphisms satisfy the ordinary algebra square.

\begin{theorem}[Native initiality]
Under guard (G) and the stated identifier convention, $(F,a)$ is initial among native $T_P$-algebras. For every $\alpha:T_P(X)\to X$ there is a unique actual native arena map $\operatorname{fold}_\alpha:F\to X$ with
\[
 \operatorname{fold}_\alpha\,a
 =\alpha\,T_P(\operatorname{fold}_\alpha).                 \tag{5}
\]
The uniqueness is on configurations, identifiers, actual row/coefficient maps and the whole unused-complete weighted grade, not on a carrier quotient. External $A$ roots are the corresponding rooted extension $T_P^A(F)$, with the same full tree-grade constructor.
\end{theorem}
\begin{proof}
Define the configuration function on every finite invocation tree by structural recursion:
\[
 h(s,(t_p))=\alpha_{\rm cfg}\bigl(s,(h(t_p))\bigr).         \tag{6}
\]
The input fibres are correct because the child at port $p$ has input $u_p$. Guard (G) says $F$'s initial configuration at $u$ is the single node $i_P(u)$ with no children. It is also the initial configuration of $T_P(X)$, regardless of $X$. Since $\alpha$ preserves initials, (6) gives $h(i_F(u))=i_X(u)$. This is the key initialization equation; it is not an assumption that the desired fold already preserves feedback.

An enabled root event maps through the identity root-event map of $T_P$ followed by $\alpha$'s native event/row map. An enabled event in a child first uses that child's recursively defined fold map and then $\alpha$'s map at the corresponding copy. Its actual raw and coefficient arrows are the composites along this finite ancestor path. Induction on the current tree gives a bijection of enabled events and exact row pushforward, since the child maps and $\alpha$ have those properties. Emitted output-port maps are defined in the same way, and are order-preserving label-preserving bijections. For a descendant identifier $((p,u)\cdot v,e)$, apply the recursively defined child identifier map to $(v,e)$ and then the global $\alpha$ map at the continuation tag $(p,u,-)$. The loop label remains available even if the same port identifier occurs under another label elsewhere. Thus these are global functions on all occurring words, not state-dependent choices; their domains are justified by the corresponding folded configurations. Port maps use the same labelled-word recursion.

For a root transition, the existing folded children stay fixed and each newly born child folds to $i_X(u_p)$ by the initial equation. Thus the transition of $T_P(X)$ creates exactly the same children as folding the new feedback configuration. Apply $\alpha$'s transition square. For an event in an existing child, use its inductive transition square and then $\alpha$'s square. This induction is on the old finite ancestor structure; it is valid even when the successor contains new nodes. The same argument for simultaneous events uses the supplied persistence coherence of $\alpha$ and the child maps, proving coherence of the actual row arrows, not just the resulting states.

Independently fold the entire grade. For a tree with root shape $c$, replace each external leaf by the $I$ option and each $U$ child by its label paired with the recursively folded child shape; apply the full grade map of $\alpha$ in (4). Its actual coefficient arrow is the tensor of all child-fold arrows and the root identity, followed by $\alpha$'s displayed coefficient arrow, with the unitors for external leaves. This defines an arity-preserving shape map and actual arrow at every finite formal tree, including all unused ones. The terminal body square follows by applying the terminal square of $\alpha$ at the root and the child induction. Content coefficients are folded through $\alpha$'s content arrows, never through event context coefficients.

These definitions prove the algebra equation (5) on every item of native data. Conversely (5) forces (6) at each root/children configuration. On an enabled event it forces the stated root map or the corresponding child-map composite, so induction along its finite ancestor path determines the raw and actual context-coefficient arrow. It similarly determines every occurring stable identifier and output-port map. The full grade equation forces the recursive shape and coefficient arrows at the root constructor of every formal tree, with no restriction to realized shapes. Thus every component of a competing native algebra morphism is equal to the displayed fold. This proves uniqueness.
\end{proof}

No initial object of the whole arena hom category was used. This is initiality for a particular independently defined endofunctor, not freeness of a general feedback doctrine. It is an open-arena universal property: a fold may identify raw histories with incompatible selected controls, so existence of this native fold does not imply that every shared policy descends through it. The separate local control relation remains necessary. The externally rooted compiler preserves its complete open interface; evaluating only initial $A$ roots may still forget unused primitive context rows, as previously distinguished.

\paragraph{Why port-only invocation words are insufficient.}
Suppose the same original port identifier $p$ emits a $U$ label in one configuration and a $V$ label in another, and the child initially enables identifier $e$. With port-only words, combined feedback uses the same global child identifier $(p,e)$ in both cases, while nested feedback distinguishes an inner-$U$ child from an outer-$V$ child. The per-configuration flattening is bijective but globally identifies these two names, so no inverse global identifier function exists. Likewise tagging the feedback word by $(p,u)$ while leaving $T_P$ continuation identifiers at $(p,e)$ would break its algebra unfolding across different $u$. The paired refinement of both constructions above is therefore material. It preserves additional already supplied label distinctions rather than quotienting any original native data.

\section{Operational nesting, empty feedback and colored-grade coherence}
Let $P:A+U+V\rightsquigarrow B+U+V$ satisfy the guard forbidding all initial $U$ and $V$ ports. The first feedback $\operatorname{Fb}_U(P)$ is an arena $A+V\rightsquigarrow B+V$. Its initial configuration is a single original initial state, so it has no initial $V$ port and the second feedback is admitted.

\begin{theorem}[Native nesting at the open control interface]
There is an operational native isomorphism
\[
 \operatorname{Fb}_V(\operatorname{Fb}_U(P))
 \cong \operatorname{Fb}_{U+V}(P),                         \tag{2}
\]
with full content-grade isomorphism
\[
 R_V(R_UC)\cong R_{U+V}C.                                 \tag{3}
\]
It preserves every enabled primitive event row and all simultaneously selectable event subsets. It therefore transports arbitrary complete selected policies, including their shared laws and zero/off-policy tables, without a scheduler change or independence assumption. The isomorphism is natural in all guarded native arena maps respecting the two loop colors. Empty-loop feedback is operationally the original arena, with $R_{\varnothing}C\cong C$. Finite repeated nestings have coherent flattenings.
\end{theorem}
\begin{proof}
A nested configuration is a finite $V$-invocation tree whose nodes are finite $U$-invocation trees of original $P$ configurations. An emitted $V$ port of a $U$-tree is identified by its stable $U$-labelled path to an original node and that node's port; its outer invocation edge also retains the actual $V$ label. Attach the outer child precisely at that labelled port and flatten the nested family. In the combined word, an inner edge $(p,u)$ carries label $\mathrm{inl}(u)$ and a $V$ edge $(p,v)$ carries $\mathrm{inr}(v)$. Splitting a combined word at its $V$-coloured pairs recovers the nested word globally. This works even when the same original port identifier occurs with different colors or different same-color labels in different configurations. The result is a finite tree of original configurations, with every internal edge colored $U$ or $V$, containing exactly one child for each emitted port of either loop color. Thus it is a configuration of the right side of (2), including for unreachable original states.

Conversely remove the $V$ edges of such a colored tree. Each remaining maximal $U$-connected component is a valid $\operatorname{Fb}_U(P)$ configuration: it retains all the $U$ children of its original nodes. The removed $V$ edges link these components into a finite outer invocation tree, attached at the corresponding stable emitted-port paths. This reconstructs the nested state. The two operations are inverse on configurations and on the global labelled identifier occurrence unions, not a quotient of states or histories.

In either configuration the enabled events are exactly the pairs consisting of an original $P$ node and one of its enabled events. The row, raw set, law and actual context coefficients are identical. Updating such an event changes the same original state, retains existing children, and initializes the same new colored children. Guarded initial states have neither color of loop port, so both routes perform exactly this finite initialization with no hidden cascade. The recursive planar order of output $B$ ports is the same flattened order. Thus the configuration bijections commute with transitions, initials, event and output-port bijections, and every structural diamond.

At termination, flattening a $V$-tree of $U$-trees gives a finite $C$-tree whose internal edges carry labels in $U+V$. Removing $V$ edges is again its unique inverse. This works for every full formal grade shape, not only realized trees. Its actual coefficient isomorphism reassociates and, where needed, permutes the same individual $D_c$ factors. Symmetric monoidal coherence gives the inverse and makes the terminal body square commute. Native maps act on all original states, stable paths and full shape factors, so both flattening routes agree nodewise. This proves naturality.

An operational isomorphism identifies all legal histories and every enabled-scope coupling port. Transport the whole selected subset and its same raw law, rather than sampling a new law or projecting to component-local histories. This proves the policy assertion. With no loop colors the configuration has just its original root and no children; the full tree grade has only single-node $C$-trees, proving the empty-loop statement. With three or more loop colors, every composite flattening yields the same original-node colored tree and event identifiers. Every content arrow is the same coherence map on its individually labelled node factors. Therefore all parenthesization and empty-color coherence diagrams commute.
\end{proof}

The nesting law is for guarded early-port recurrence with changing full grades. It is not ordinary fixed-grade yanking or a claim that arbitrary effectful sliding is invertible. In particular it does not rely on contracting an unfinished arena into an atomic complete-execution row.

\section{Chosen controls, native cuts and terminal event clock}
At every legal feedback raw past, including all zero and off-policy pasts, a deterministic scheduler selects a nonempty finite subset of enabled pairs $(v,e)$. A supplied normalized joint law has the original event marginal at each selected event conditional on the complete paid past. The selected scope is atomic. Complete controls retain all these tables; evaluated histories follow the scheduler choices and retain every zero draw. Events in a selected scope can belong to different invocations or to an unfinished parent and an already started child.

Normalized finite prefix laws follow by induction on paid scopes. Each prefix coefficient is the tensor of its executed event coefficients. Raw last-scope truncation has exact fibre pushforward and an actual context augmentation-deletion map. Its target is the retained operational configuration/history interface, not an alleged polynomial $C$ shape for an unfinished invocation. A normalized $R_UC$-graded prefix process is not constructed: terminal-only shape data do not authorize such a body. A separate configuration-annotated polynomial prefix completion would require its own grade and truncation proof.

Under a native arena map, require the selected scheduler subset to map to the target subset and the product raw map to push its conditional law exactly, at every source legal past. Finite induction then gives actual native prefix/context transport and exact terminal raw pushforward; terminality is reflected by the enabled-event bijections. The terminal coefficient arrow tensors all event maps and the global grade map is the full $R_u\gamma$. Bounded stopping cuts must be defined on all legal pasts and occur between entire atomic scopes. Retain the complete paid past, all invocation configurations, pending ports and complete suffix tables, including off-policy reserve. Concatenating the raw scope sequences and partitioning their finite event tensors proves actual context cut/assembly and its native coherence. It does not turn a marked invocation-tree ideal into an admissible cut if it splits a paid scope or discards an unfinished invocation.

For the terminal event-count clock assume countable $A,B,U$ so the reviewed countable-endpoint trace applies; finite branching still makes all reachable finite raw histories countable without any countability assumption on the full content grade. The clock index is $n\in\mathbb N_0$, explicitly including zero. At row $(a,n)$, execute at most $n$ nonempty paid scopes, absorb terminals, and use the explicitly additional typed restriction to terminal histories with exactly $n$ paid events. Forget the complete event schedule in the terminal body but retain its raw history and tensor coefficient. Regrade into the entire $R_UC$ and return clock $n$.

\begin{theorem}[Chosen-control terminal trace compiler]
This clock is admitted by the countable subprobability matching trace and its trace is the chosen-control terminating feedback evaluator, with every terminal zero witness and identity full $R_UC$ grade. It is natural for the stated conditional-compatible native comparisons and for tensor-to-product comparison on the same event factors. No initial-output cemetery label or mid-invocation terminal shape is introduced.
\end{theorem}
\begin{proof}
A terminal history with $n$ events has at most $n$ selected nonempty scopes, so it appears in that finite unrolling. Absorption records it once; the filter selects precisely its event-count layer. For $n=0$, an initially terminal feedback configuration is already included. Guard (G) ensures the initial feedback configuration itself is a finite state even in that case. Each layer has mass at most one, and all layers partition the finite terminal histories. Finite-horizon conservation therefore gives total layer mass at most one, the exact trace admission test. Matching the clock collects each terminal history once, without modifying its actual coefficient, body or full grade. Native maps satisfying the local equations preserve event count and transport the finite prefix fibre equations and terminality, hence restrict to this layer and trace. The tensor-to-product maps preserve raw histories and full grades and commute with all finite event deletions and partitions, so their restriction and trace squares commute as well.
\end{proof}

The operational nesting isomorphism (2) preserves the same primitive events and selected scopes, hence their exact event count. It consequently commutes with every finite clock unrolling, terminal restriction and countable trace, with full grade map (3), for all transported selected policies. This is stronger than merely comparing the independent terminal laws.

The restriction is not an old exact-pushforward map; regrading the selected row occurs in the subprobability category. All nonterminal prefixes remain in the separate operational construction. Nontermination mass is not assigned to an invented divergence witness.

\section{Fixed-full-grade recurrent prefix compatibility}
There is now an actual graded prefix interface without inventing a completed shape for an unfinished invocation. Use the partial native category of \emph{Native recurrent prefixes with partial terminal bodies}. An object is a normalized native context $P_n$, the full collection $R_UC$, a terminal subdomain $T_n\subseteq e_AP_n$, and a body $b_n:T_n\to F_{R_UC}B$. An arrow $(r,\rho)$ satisfies
\[
 r^{-1}(T')\subseteq T,\qquad F_\rho b=b'r\text{ on }r^{-1}(T').       \tag{7}
\]
Thus truncation may send a newly terminal source to an active predecessor, but cannot send an active source to a terminal predecessor. On full-domain objects this is exactly the original native comma category; no new equality of old arrows is introduced. No content augmentation is assumed.

\begin{theorem}[Observed partial-body tower of guarded feedback]
Every selected guarded feedback arena has a normalized native inverse tower at the \emph{same entire} grade $R_UC$, with terminal-only bodies, actual tensor context truncations, and retained paid-history/configuration observations. Conditional-compatible native arena maps induce strict-domain native program comparisons of these towers. Operational colored nesting gives an isomorphism of these observed strict towers, with its full colored-grade isomorphism, for all transported selected policies. Their absorbing first-hit reconstruction is the same terminating process as the event-count clock.
\end{theorem}
\begin{proof}
At scope horizon $n$, take precisely the scheduler-compatible histories stopped at that horizon, absorbing earlier terminals as themselves. Local finiteness and guard (G) make every finite-horizon raw row finite. Its weight is its product of normalized scope probabilities, and its coefficient is the tensor of its executed event factors. The finite-horizon mass induction proves normalization. Let $T_n$ be exactly the terminal histories and use the already constructed full $R_UC$ terminal body there only.

Raw truncation from $n+1$ to $n$ has exact fibre pushforward: above an active prefix sum its full next-scope law, while an old terminal prefix has a unique absorbed copy. Its actual context arrow augments only the new event factors. If its target prefix is terminal then its unique source continuation is terminal with unchanged body. This proves (7) with identity on the whole $R_UC$ and shows precisely why equality of domains is not required for truncation. Content coefficients are never deleted; the same full grade is carried at every horizon.

Additionally retain the native contextual observation $o_n:e_AP_n\to V_n$ into the set of complete paid histories at that horizon, decoding all invocation configurations, stable early ports and pending inputs. It commutes with literal history truncation. The external input remains $A$, and its paid context weights and coefficients are not replaced by an uncosted history input. The observed tower retains the complete policy as an index; neither off-policy tables nor zero-history choices are reconstructed from the scalar weights. A bare partial-comma reduct without $o_n$ and the arena index would not preserve active early-output distinctions.

Under the stated policy-compatible native maps, finite-step pushforward and tensor maps are already proved. Enabled-event bijections reflect terminality, giving equality $f_n^{-1}(T'_n)=T_n$ at every horizon, in addition to the full $R_u\gamma$ body square and history-observation square. They therefore are strict-domain program comparisons. The nesting isomorphism identifies the same scope histories and observations at each horizon and the same individual event factors. Its actual full grade map (3) then gives the asserted levelwise native tower isomorphism, not a reindexed pro-object comparison.

Old terminal atoms have unique absorbed copies, unchanged weights and coefficients. Take their disjoint first-hit layers, including the initial layer $T_0$. Each finite terminal history occurs once and their total mass is at most one. Strict-domain program maps preserve these layers and exact fibres, yielding their native terminating map. The event-count clock partitions the same retained terminal histories by a different finite integer. After assembly both raw sets have the identity identification on complete terminal histories, with identical tensor coefficients, body and full grade. No map between the unclosed step and event clocks is inferred.
\end{proof}

This supplies E4 prefix compatibility in a changed \emph{partial-body} interface while retaining original full content maps. It does not claim a total-$R_UC$ polynomial body at active prefixes or functorial terminal reconstruction for every lax arrow satisfying (7). The latter would be false: an identity context map from a constant terminal-domain tower to a constant empty-domain tower satisfies (7), yet their reconstructed terminal masses are one and zero. Actual program maps have the stronger domain reflection proved above.

\section{Independent invocation normal form}
Let $f=\operatorname{Eval}_{\rm ind}(P):A+U\to F_C(B+U)$ have as raw witnesses the complete finite local event traces modulo exchanges of adjacent simultaneously enabled events, retaining all raw outcomes including zeros. Its weight is the product of primitive event weights; its context coefficient is their actual tensor, and its body is the terminal $P$ body. Finite branching gives a countable raw set. Its total mass is at most one, not necessarily one: $P$ may have an infinite run. Thus $f$ is a countable \emph{subprobability} process.

The invocation-level independent recursive evaluator of $f$ is defined directly by independent subprobability bind on finite invocation trees. Its weight is the product of the $f$-witness weights and its context coefficient the tensor of their coefficients. A product of subprobability rows is not generally a coupling with those original marginals, because its marginal includes the total masses of the other factors. No normalized coupling condition is asserted for this macro evaluator. Finite-generation mass is at most one by the subprobability bind calculation, and countable summation gives its terminating subprobability law.

\begin{theorem}[Native independent feedback/iteration comparison]
For every deterministic nonstalling independent primitive-event scheduler, the terminal feedback process is natively isomorphic to this invocation-level independent subprobability evaluator of $f$, at identity full grade $R_UC$. The comparison preserves actual tensor context coefficients and all native guarded arena maps. It concerns finite terminal behavior; it does not identify recurrent prefixes or arbitrary shared controls.
\end{theorem}
\begin{proof}
First pass from a selected independent finite terminal history to its complete event trace under the compiled diamonds. Conversely a deterministic scheduler reconstructs one history from a complete finite trace. Any currently enabled event must appear in a terminal extension, since persistence prevents its removal and a terminal state has no enabled events. It can be exchanged leftward across the intervening events, keeping its literally stable row and raw outcome. Repeat for the prescribed subset and the remaining finite trace. This terminates within the trace's finite event count and does not need a global bound or fairness. The weight is the product over primitive events and the coefficient map is the actual tensor permutation, including zero outcomes.

Restrict a complete feedback trace to each invocation. Other invocations never change its local configuration, so this is a complete local $P$ trace beginning at its designated initial state. Its terminal emitted $U$ ports identify exactly the child invocations, with their inputs and stable identifiers; its terminal full shape fixes their final ordered positions. These data form precisely a finite raw invocation tree for $f$.

Conversely, given such an invocation tree, run the complete root trace first and then its child invocation traces in their final port order, recursively. Each child has been initialized by the time it is used, so this is valid. Different representative orders inside a local trace are related by its diamonds. An arbitrary original interleaving can be converted to this invocationwise order: move each remaining parent event left across earlier child events after its own local predecessors have executed. They commute because an existing child is unchanged by later parent events. Then apply the same argument inside each child and between disjoint children. Thus the two constructions are inverse on complete trace classes, not just surjective encodings.

Grouping the product of primitive weights by invocation gives exactly the product of the $f$-witness weights. Grouping the actual augmented event tensor by invocation gives the context coefficient of the macro tree via canonical associators and symmetries. The full body renders the same invocation tree, with its independent per-invocation content factors and final ordered leaves, so the global map is identity on all $R_UC$. These give native inverse arrows.

A native arena map preserves event/port diamonds, complete local trace classes, terminal shapes and stable child addresses. Its event maps preserve event count. For exact local trace pushforward, choose a target independent scheduler and pull each selected subset back through the native enabled-event bijection at every source history. Independent product laws push to the target products by the primitive fibre equations. Finite-step induction now applies to genuinely matched schedules, and terminality is reflected by the enabled-event bijections. Countable summation over terminal lengths gives exact terminal pushforward. Scheduler invariance then identifies these rows with any other chosen independent evaluators; no incompatible-prefix map is invented. The induced invocation-tree pushforward follows by independent raw fibre sums at each finite target tree and nonnegative Fubini. On coefficients both routes tensor the same actual event maps; on content both use the full $R_u\gamma$. Hence the isomorphism is native-natural. This does not claim a faithful evaluator on all unreachable primitive configurations or a compatible prefix map between arbitrary different schedulers.
\end{proof}

For an atomic original process embedded as a one-paid-event arena, the guard holds automatically. The local trace is exactly the original raw witness and its coefficient, so the theorem recovers original independent recursive raw trees and actual tensor maps at the full $R_UC$ grade. General open arena implementations have additional early-event control ports; this recovery does not contract those ports for shared policies.

\section{The cost clock is the common trace interface}
Assign each macro raw witness $w$ its cost $c(w)$, the number of primitive events in its complete $P$ trace. A zero-cost witness starts and ends at the initial configuration, so (G) implies its body has no $U$ output. It may have nontrivial $B$ outputs and content coefficient; its context coefficient is the tensor unit. Every invocation with a recursive child therefore has positive cost.

\begin{proposition}[Finite cost bound]
A terminal invocation tree of total cost $n$ has depth at most $n+1$, even when some invocations have cost zero. Hence its complete cost-$n$ layer is obtained by $n+1$ generations of independent subprobability macro unrolling followed by typed restriction to terminal trees of total cost exactly $n$.
\end{proposition}
\begin{proof}
Along a root-to-leaf invocation path every node except possibly the last has a $U$ child and thus cost at least one. A path with $d$ nodes therefore contributes at least $d-1$ to the nonnegative total cost of the whole tree, so $d\leq n+1$. Breadth-first macro generation reaches every node within that many generations. Terminal absorption avoids duplicates. Restricting by total cost selects exactly the required layer. No bound on the total number of zero-cost leaf invocations is needed; each individual raw tree is finite, and their row set is countable. For $n=0$, one generation suffices for an initially terminal $B$-only root.
\end{proof}

\begin{theorem}[Fixed-cost native trace square]
The event-count clock of independent feedback and the macro cost clock have native isomorphic rows at every fixed clock value $n$, identity full $R_UC$ grade, and the same countable traced terminating process. The square is natural in guarded native arena maps and commutes with the directed tensor-to-product comparison. It requires neither equality of scheduler step counts nor a guessed map from an invocation count to a primitive event count.
\end{theorem}
\begin{proof}
The independent normal-form bijection preserves total primitive event count, which is exactly the sum of the macro witness costs. Its raw fibre equation and actual tensor isomorphism therefore restrict within each fixed $n$. The event row is constructed by $n$ paid-scope steps; the macro row is constructed by $n+1$ generations by the proposition. Both restrictions select the same finite terminal trace data after their respective canonical normal forms. Native arena maps preserve event count and local trace costs, so the naturality square restricts too. Countable trace uses the unchanged coefficient arrows and identity full grade. Tensor-to-product comparison commutes with the partition of all primitive events into invocation blocks, including zero-event blocks, by its established empty-block and finite-partition equations. Thus its two routes agree as actual arrows.
\end{proof}

The product comparison in this square has an explicit two-stage target. If $E_v=\bigotimes_{e\in v}K_e$ is an invocation's tensor coefficient, then
\[
 \bigotimes_v E_v
 \xrightarrow{\lambda_{\rm inv}}\prod_v E_v
 \xrightarrow{\prod_v\lambda_{\rm events(v)}}
 \prod_v\prod_{e\in v}K_e\cong\prod_eK_e
\]
is the eventwise product comparison. A product of invocation tensors is not identified with the product of individual event coefficients: the second displayed arrow is additional and can have a kernel. The partition calculus proves equality only after that arrow, including empty event blocks.

This is a trace-generated assembly theorem with its additional generators displayed: guarded early-port initialization, finite operational unrolling, explicit typed restriction, and countable matching. It is not the assertion that general instantaneous trace selects a least fixed point. The initial-algebra fact (1) is about finite current configurations, while the chosen-control terminating law and the independent comparison are separate semantic statements.

\section{Strict shared-feedback extension and residuals}
Here is a guarded finite implementation separating early feedback from invocation-atomic recursion. At an external root, execute a fair event $X_1$ and emit one $U$ port immediately. Then execute a second fair event $X_2$ and emit a $B$ port labelled by $X_2$; the root invocation's full terminal shape is binary, ordered as the earlier $U$ port followed by this $B$ port. At the loop input, the arena executes one fair event $Y$ and emits one $B$ port labelled by $Y$, with unary terminal shape. Every initial state has no $U$ port. The complete root macro row is the independent fair pair $(X_1,X_2)$, while the child macro row is fair $Y$.

In the feedback arena, after $X_1$ the unfinished parent's $X_2$ event and the child's $Y$ event are simultaneously enabled and have their unchanged fair rows. Choose the diagonal coupling for that scope. It is admitted conditional on the entire earlier past, and the two final ordered output bits $(Y,X_2)$ agree surely. All off-diagonal raw scope tuples remain with weight zero. Give each event the nontrivial augmented coefficient $k^2\to k$; the context tensor retains all three factors, while the full content grade remains all $R_UC$.

In invocation-atomic causal recursion of the same macro rows, the entire parent pair $(X_1,X_2)$ is sampled before the child is enabled. Conditional on that past the single child's primitive row is still fair. Consequently $\Pr(Y=X_2)=1/2$ for every admitted invocation-atomic conditional policy. It cannot produce the event-level diagonal law without changing that required conditional marginal. This is a strict law-level extension at the specified interface, not a claim that every conceivable macro formalism is incapable of representing it.

\paragraph{Recovered unbounded shared branching interpretation.}
The earlier balanced-tree shared recurrence is an instance of this actual compiler, not a new standalone theorem. Let $P$ be a one-event arena with one loop state and two equiprobable outcomes: a unary terminal $B$ shape or a binary shape returning $U,U$. Give every event witness the nontrivial augmented coefficient $k^2\to k$, and retain independent weighted content coefficients at both full shapes. At each frontier choose a fresh fair bit and use that same branch/terminate choice at every enabled event; mixed raw tuples remain present with probability zero. Initial configurations emit no port, so the guard holds.

A positive terminal history consists of $d-1$ complete binary branching generations followed by one all-terminal generation. It has probability $2^{-d}$, $2^{d-1}$ settled leaves and $2^d-1$ executed events, with their full nontrivial context tensor. Thus termination probability is one, while expected event count is infinite: the partial sum through depth $N$ is $N-1+2^{-N}$. The full $R_UC$ nevertheless retains every unbalanced and mixed tree and its content coefficient. Three-leaf raw trees have mass zero for this shared policy; under independent event choices their two planar forms have total mass $2\cdot2^{-5}=1/16$. The compiler's observed prefix tower, native initiality, nesting and event-cost trace all apply to this one implementation. Neither a support-only tree grade nor an infinite coefficient tensor is required.

The guard is not a theorem about unrestricted feedback. Without it, an initial $U$ port can demand another identical initial copy before any paid event, creating an infinite silent initialization cascade outside this finite-configuration compiler. A finite or well-founded silent-closure interface would be a different repaired alternative. Likewise, fairness/liveness of individual pending calls is separate from finite-terminal independent semantics, and full infinite-execution coefficients or laws are not constructed. A completed invocation shape may be unknown during its early outputs, so no full-$C$ partial-prefix grammar is inferred; all those configurations and ports remain explicit operational data. Terminality in this arena means an empty enabled set with a supplied terminal body; it does not automatically mean successful termination in an application. Refusal, deadlock and distinct completion classes require explicit output labels or additional operational outcome structure if they are to be distinguished. These residuals delimit, rather than erase, the direct native relationship established here.
\end{document}
